\label{sec:empirical}

The implementation slice for M.3 is now present in \textsc{bisect}. The
\texttt{acs-housing} fetch path derives the four housing-character columns, and
the \texttt{housing-character} weight mode applies the cosine-similarity blend
to graph edges. This section separates implemented evidence from future
empirical claims: the formula and edge-weight mechanics are package verified,
while state-level outcome results remain future work until full NC-14 runs are
archived.

\subsection{Target State and Configuration}

The primary empirical target is North Carolina, 14 congressional districts,
2020 census vintage.
North Carolina provides an ideal test case because it contains all three
prototypical housing character communities within a single state:
\begin{itemize}
  \item \textbf{Historic neighborhoods.}
    Raleigh's Cameron Park, Five Points, and Oakwood;
    Durham's Trinity Park and Watts-Hillandale;
    Charlotte's Dilworth and Myers Park.
    These neighborhoods have median year built predominantly before 1950,
    high owner-occupancy, and predominantly single-family or attached rowhouse
    stock.
  \item \textbf{Postwar single-family suburbs.}
    Cary, Apex, Morrisville, and the Research Triangle Park periphery.
    Median year built 1970--2000, high owner-occupancy, predominantly detached
    single-family.
  \item \textbf{Apartment corridors.}
    Charlotte's South End and South Park; Raleigh's North Hills and Glenwood
    South; Durham's downtown apartment district.
    Predominantly large multifamily, high rental fraction, recent construction.
\end{itemize}

The baseline comparison is \texttt{--weights-override geographic} on the same
NC-14 plan configuration (convergence search, $T = 600$, balance tolerance
$= 0.5\%$).

\subsection{Primary Empirical Hypothesis}

\paragraph{H1: Within-district housing homogeneity improves.}
The M-track primary metric for M.3 is the mean within-district standard
deviation of \texttt{pct\_single\_family} across all 14 NC districts.
Formally:
\begin{equation}
  \overline{\sigma}_{\mathrm{SF}}
  = \frac{1}{14} \sum_{i=1}^{14}
    \mathrm{std}\bigl\{\mathrm{SF}(t) : t \in D_i\bigr\}.
  \label{eq:within-sf}
\end{equation}
We predict that $\overline{\sigma}_{\mathrm{SF}}$ under housing-character weights
will be at least 10\% lower than under geographic weights. This is a future
empirical claim, not established by the smoke package.

\paragraph{H2: Historic neighborhood co-assignment.}
Define the set of ``historic tracts'' as those with $\mathrm{VT}(t) > 0.64$
(median year built before 1960) in North Carolina.
We predict that the fraction of historic tract pairs that are co-assigned to
the same district under housing-character weights will be at least 60\%,
versus at most 40\% under geographic weights. The threshold remains a
pre-registered target for the future state run.
This tests whether the vintage signal causes the historic neighborhoods of
Raleigh and Durham to cluster together rather than being split across suburban
and apartment-corridor districts.

\paragraph{H3: Apartment corridor isolation.}
Define ``apartment corridor tracts'' as those with $\mathrm{MF}(t) > 0.50$
(majority large-multifamily units).
We predict that the mean within-district $\mathrm{MF}$ variance under
housing-character weights is lower than under geographic weights,
indicating that apartment corridor tracts are grouped together rather than
scattered across single-family districts.

\subsection{Package-Verified Mechanics}

The smoke package
\texttt{docs/examples/m-community-character-evidence-packages/M.3+acs-housing-smoke}
contains two synthetic tract rows and one edge. The verifier recomputes:
\begin{itemize}
  \item \texttt{pct\_single\_family} from B25024\_002E and B25024\_003E over
    B25024\_001E;
  \item \texttt{pct\_multifamily} from B25024\_007E--B25024\_009E over
    B25024\_001E;
  \item \texttt{pct\_owner} from B25003\_002E over B25003\_001E;
  \item \texttt{housing\_vintage} from B25035\_001E with the 1940--2020 clamp;
  \item the final edge weight
    $w_h = w_{\mathrm{geo}}(\alpha + (1-\alpha)\mathrm{sim}_h)$.
\end{itemize}
The package is deliberately synthetic. It validates implementation mechanics
and tamper rejection; it does not claim an empirical compactness,
proportionality, or community-preservation result for North Carolina.

\subsection{Mill Hill Clustering --- Qualitative Prediction}

North Carolina has documented historic mill communities: the village neighborhoods
built by textile mills in the 1910s--1940s around Durham, Burlington, and
Kannapolis.
Mill Hill neighborhoods share a distinctive housing character: pre-1950
construction, attached or detached single-family worker housing, moderate
to high owner-occupancy.
They are geographically proximate to industrial parcels (former mill sites)
but have high housing character similarity to each other across municipal
boundaries.
We predict that under housing-character weights, Durham's Mill Grove and
West End, Burlington's Edgewood, and similar mill villages will show elevated
co-assignment rates compared to geographic weights.

\subsection{Implementation Cost Confirmation}

A secondary empirical goal is to confirm the zero-cost implementation claim:
that adding B25024, B25003, and B25035 to the ACS fetch requires no new
download infrastructure, no new API keys, and no changes to the adjacency
graph pipeline.
This has been verified for the formula and graph-weight layer by:
\begin{enumerate}
  \item extending the ACS housing CSV with the four new columns;
  \item adding formula tests for denominator defaults and vintage clamping;
  \item adding a hash-bound fixture that confirms the housing-character weight
    mode produces a positive blended edge weight.
\end{enumerate}
The full all-tract North Carolina fetch and plan comparison remain future
empirical work.

\subsection{Comparison Methodology}

All comparisons between housing-character and geographic weights will use
the same seed (or convergence ensemble with $T = 600$) to isolate the effect
of the weight change.
The proportionality gap (D-seat share minus D-vote share) will be reported
for cross-paper comparability with B-track results, with the note that
proportionality is not the primary M.3 target metric.
